\figwide[H]{fig01_teaser}{The paper in one picture. One real setup (game~4 of the Worlds~2025 final, first baron, 60\,s after the spawn). \emph{The map:} walkable graph (white: walls) coloured by which team reaches each cell first, with each formation drawn as the blocks of a set partition. \emph{The setup:} each team's walk on the partition lattice $\Pi_5$ (up = fewer blocks) and its players at the pit. \emph{The choice:} $k$-arrival times $\tau^{(k)}$ and the LB-QRE action $(k,m,\tau)$ at the commitment time, with win probability $W$, fight probability $\hat p$ and the objective's value.}{fig:teaser}

\section{Introduction}
\label{sec:intro}

Neutral objectives in \emph{League of Legends} (LoL) appear on a public clock.
The first dragon spawns at 5:00, the void grubs at 8:00, the Rift Herald at 15:00, Atakhan at 20:00 and Baron Nashor at 25:00, and both teams know these times to the second.
In the minute before each spawn the two teams move their ten players into position: some walk to the pit, some keep pushing a lane, some wait in the jungle in between.
This setup decides who fights for the objective, with how many players, and whether the other side trades it for towers elsewhere.
Professional coaches plan it in detail. Esports analytics has measured almost none of it.

The closest work is TAGNet \citep{zhang2026tagnet}, which studies how defenders set up after the freeze time in professional \emph{Counter-Strike~2}.
TAGNet embeds the five defenders as a complete graph with a graph attention network \citep{velickovic2018graph} and an untrained recurrent layer, clusters the 64-dimensional embeddings with $K$-Means at a fixed $K=6$, and reads the win rate of each cluster as the quality of a tactic.
Three steps of that pipeline make the conclusions fragile.
The graph carries no map: straight-line distances ignore walls.
The context variables (economy, score) are fed into the representation and later `discovered' to be associated with it.
And comparisons of trailing and leading teams average over teams rather than over rounds.

We take a different view. A setup is an \emph{allocation on a walled map}.
Distances are geodesic on a navigation graph built from Riot's own position data, so a formation is the set of connected components of a proximity graph, that is, a set partition of the five roles, and the formations form the partition lattice $\Pi_5$ with \V{latticeN} elements.
Map control is the geodesic Voronoi split of that graph, and readiness is a vector of $k$-arrival times.
The decision is how many players to commit to the pit, how many to a second front, and whether to be there before the spawn.
Fight outcomes follow a Lanchester-type law \citep{lanchester1916aircraft} written as a Tullock contest \citep{tullock1980efficient}, the second front is a Blotto-type contest \citep{roberson2006colonel}, and every outcome is priced in win probability by a surface $W$ fitted on game results only.
\Cref{fig:teaser} shows all of these objects on one real baron setup from the Worlds~2025 final.

The first finding is about behaviour. Across all \V{nGames} games of MSI and Worlds~2025, trailing teams leave the pit empty far more often than leading teams once the objective is up.
At the first baron the gap reaches \V{vCBaronFortyfive}\pp{} \V{vCBaronFortyfiveCI} at $s{+}45$\,s, at Atakhan \V{vCAtakhanFortyfive}\pp{} \V{vCAtakhanFortyfiveCI}, and at respawned dragons \V{vCDragonTwoFifteen}\pp{} \V{vCDragonTwoFifteenCI} (\Cref{fig:curves,fig:contrasts}).
Every headline contrast keeps its sign at three pit radii and survives false-discovery control.
The same contrast on Riot's exact 60-second frames gives \V{rBaronEmpty}\pp{} \V{rBaronEmptyCI} at the first baron.
The dragon gradient of our kill-anchored pilot (\V{pilotDragonTrail}\% of trailing vs \V{pilotDragonLead}\% of leading non-takers contested) disappears once spawns rather than kills anchor the window: it was an artefact of selecting on the outcome.

The second finding is about structure. HASSE, a sticky hidden Markov model whose emissions live on $\Pi_5$, describes each setup as a walk on the Hasse diagram with a pattern-specific centre path, rigidity and inertia.
Held-out likelihood improves up to $K=8$ with steps that shrink to \V{hvStepSevenEight} nats per setup, and six patterns carry the structure (\Cref{fig:patterns}).
Two of them, used almost only at early river objectives, stay spread out; two gather four players around the pit at the spawn and dominate the late game.
Whole-model series-bootstrap refits of HASSE agree at ARI \V{hsBootKSix} (\V{hsBootKSixN} series-bootstrap resamples), comparable to re-running only TAGNet's $K$-Means step on fixed embeddings (\V{tgBootAri}); five independent HASSE fits, each the best of ten restarts, agree at \V{hsCrossKSix}, whereas single TAGNet runs on the same V-layer tracks agree at \V{tgSeedAri} across seeds and at \V{tgCountAri} with TAGNet's own counting baseline (\Cref{fig:stability}).

The third finding is about incentives. LB-QRE, a heterogeneous logit quantal response equilibrium \citep{mckelvey1995quantal} of the allocation game estimated by nested pseudo-likelihood \citep{aguirregabiria2007sequential}, finds that teams respond to fight values, travel and the cost of each committed player, and barely to the objective's own value: its weight is $\nu=\V{qNu}$ \V{qNuCI}, and fixing $\nu=1$ (pure win-probability maximisation) costs \V{qWinmaxPoint} log-likelihood points (likelihood ratio about \V{qWinmaxLR}; bootstrap mean \V{qWinmax} \V{qWinmaxCI}).
The equilibrium reproduces about a quarter of the observed trailing-minus-leading commitment gap (\V{gapModel} vs \V{gapObs}\pp), but almost all of that comes from state variables (alive counts, the gold term, travel and feasibility): with every strategic channel off the gap is \V{gapNone}\pp, so the four channels add \V{gapChannels}\pp{} net.
Within the model, the curvature of $W$ pushes trailing teams towards contesting (\V{shC}\pp) and team behaviour pulls them away (\V{shB}\pp), while the Lanchester and second-front channels are small (\Cref{fig:qre}).
Fights punish the smaller side hard: the larger side loses the next player far less often than the square law predicts (\V{sSevenFiveThree} vs \V{sSevenFiveThreeSq} at 5v3).

These results rest on a measurement layer that we validate end to end: a broadcast-clock reader aligns every video to game time, the minimap is registered to the navigation graph, and a per-game tracker is scored against Riot positions at held-out minutes (median error \V{trkMed} map units on a map 14,820 units wide).
In summary, we contribute
(i)~a theory-defined measurement of objective setups (navigation graph, geodesic Voronoi control, $k$-arrival order statistics, partitions on $\Pi_5$) with validated extraction from broadcast video;
(ii)~HASSE-QRE, a two-layer model whose discovery layer lives on the partition lattice and whose explanation layer is an estimable allocation game in win-probability units;
(iii)~the findings above, including the negative ones, with an estimator audit that shows how the choice of estimator can flip significance; and
(iv)~an open dataset of \V{nSpawns} spawns, \V{nSetups} team-setups and 2-fps tracks, with an appendix that turns narratives about T1 and Faker into measured quantities.
